\documentclass[master, fontsize=12pt]{diploma}

\usepackage[firstConsultant]{review}

\student{Инютин Максим Андреевич}
\faculty{№ 8 <<Компьютерные науки и прикладная математика>>}
\department{Образовательный центр Института № 8}
\speciality{01.04.02 Прикладная математика и информатика}
\profile{Продуктовая разработка и разработка программных систем}
\group{М8О-206СВ-24}

\theme{Разработка алгоритма лидарной одометрии беспилотного автомобиля, учитывающего перекрытие поля зрения}
% Дата. Оставляем пустое место для дня
\workDate{\uline{\hspace{24pt}} мая \the\year\ года}

\supervisor{Ухов Пётр Александрович}
\supervisorCredentials{кандидат технических наук, доцент, преподаватель кафедры 806}

\firstConsultant{Евсеев Дмитрий Сергеевич}
\firstConsultantCredentials{доктор физико-математических наук, профессор, профессор кафедры 806}

% \secondConsultant{Алешина Елена Константиновна}
% \secondConsultantCredentials{доктор физико-математических наук, профессор, профессор кафедры 806}

% \reviewer{Свиридов Михаил Артемьевич}
% \reviewerCredentials{доктор физико-математических наук, профессор}

\headOfDepartmentTitle{И. о. директора Образовательного центра}
\headOfDepartment{Крылов Сергей Сергеевич}



\begin{document}
\fillReviewHeader

В выпускной квалификационной работе рассмотрена актуальная задача повышения устойчивости лидарной одометрии беспилотного автомобиля при частичном перекрытии поля зрения крупными динамическими объектами. Автор последовательно связывает прикладную постановку задачи с алгоритмической частью: анализирует ограничения регистрации облаков точек, описывает используемый набор данных Boreas-RT, формализует сценарий обгона грузовиком и реализует экспериментальный конвейер для проверки качества оценки движения.

С содержательной точки зрения работа демонстрирует хорошее владение методами трёхмерной геометрии и статистической обработки результатов. В предложенном решении используются двойная вокселизация облаков точек, взвешенный ICP с функцией Гемана---Макклюра, оценка евклидова преобразования по алгоритму Кабша, модель постоянной скорости и фильтр Калмана, стабилизирующий начальное приближение для последующих запусков ICP. Особенно важно, что алгоритм не ограничивается формальным описанием: он реализован программно и проверен на серии вычислительных экспериментов с различными дорожными сценами.

Полученные результаты показывают практическую значимость выполненной работы. На коротких сценах с моделируемым перекрытием поля зрения применение фильтра Калмана уменьшает влияние редких сильных деградаций оценки движения, при этом на длинных проездах естественная динамика автомобиля не подавляется. Такой результат важен для прикладных систем автономного вождения, где даже короткий участок ухудшенной наблюдаемости может стать критическим для локализации и планирования.

Заключение: выпускная квалификационная работа выполнена на хорошем научно-техническом уровне, соответствует требованиям, предъявляемым к магистерским диссертациям, а её автор заслуживает присвоения квалификации магистра по направлению подготовки 01.04.02 <<Прикладная математика и информатика>>.

\makeReviewSignature
\end{document}